% GENERATED FILE. Edit canonical lessons, then run npm run generate:books.
% canonical-source-hash: fnv1a64:b21155afc7e91fc5
% canonical-lessons: GU-C193-disa, GU-C193-dabu, GU-C193-jamnu, GU-C193-sidhu, GU-C193-valvu-turn

\chapter{Which way to go}
\label{ch:gu-which-way-to-go}
\hlchaptermodality{193}

\begin{chapteropening}
\textbf{By the end of this chapter:} \emph{I can read and give directions: left, right and straight on.}

Complete the last lesson of chapter 193: I can read and give directions: left, right and straight on.
\end{chapteropening}

\section[\texorpdfstring{\gu{દિશા} (diśā)}{દિશા (diśā)}]{\gu{દિશા} (diśā) — direction}

\label{lesson:GU-C193-disa}



\begin{quote}
Before the new one: say the Gujarati for bill, then the Gujarati for to waste.
\end{quote}

\subsection*{You'll want to know: \gu{દિશા}}
\textbf{\gu{દિશા}} — \emph{diśā} — \textquotedblleft{}direction\textquotedblright{}.

\begin{quote}
\textbf{\gu{સ્ટેશન} \gu{આ} \gu{દિશામાં} \gu{છે}.} — \emph{sṭeśan ā diśāmā̃ che} — \textquotedblleft{}The station is in this direction.\textquotedblright{}
\end{quote}

A feminine noun.

\begin{cousinweb}
\textbf{\gu{દિશા}} (\emph{diśā}) is Sanskrit \emph{diśā}, \textquotedblleft{}direction, region\textquotedblright{}, a form of \emph{díś}, taken over whole as a learned loan, a tatsama (CDIAL 6339).
\end{cousinweb}

\begin{grammarlens}[title={What you've built}]
One more word for everyday life at A2.
\end{grammarlens}

\subsection*{Your turn}
\begin{itemize}
\raggedright
  \item \emph{Say it:} \emph{diśā}
  \item \emph{Say it:} \emph{diśā}, once more
  \item \emph{Say it:} \emph{sṭeśan ā diśāmā̃ che} — The station is in this direction.
  \item \emph{Recall:} say the Gujarati for bill, then the Gujarati for to waste, then say \emph{diśā} again
\end{itemize}

\begin{tcolorbox}[breakable,colback=teal!4,colframe=teal!35!black,title={Before you move on}]
What does \gu{દિશા} mean? (\textquotedblleft{}Direction\textquotedblright{}.) Say it once more.
\end{tcolorbox}
\section[\texorpdfstring{\gu{ડાબું} (ḍābũ)}{ડાબું (ḍābũ)}]{\gu{ડાબું} (ḍābũ) — left, left-hand}

\label{lesson:GU-C193-dabu}



\begin{quote}
Before the new one: say the Gujarati for to waste, then the Gujarati for direction.
\end{quote}

\subsection*{You'll want to know: \gu{ડાબું}}
\textbf{\gu{ડાબું}} — \emph{ḍābũ} — \textquotedblleft{}left, left-hand\textquotedblright{}.

\begin{quote}
\textbf{\gu{મંદિર} \gu{ડાબે} \gu{છે}.} — \emph{mandir ḍābe che} — \textquotedblleft{}The temple is on the left.\textquotedblright{}
\end{quote}

\textbf{\gu{ડાબે}}, \emph{ḍābe}, 'on the left'.

\begin{cousinweb}
\textbf{\gu{ડાબું}} (\emph{ḍābũ}) is listed by CDIAL, glossed \textquotedblleft{}left\textquotedblright{}, under \emph{*ḍavva}, \textquotedblleft{}left-hand, crosswise\textquotedblright{}, and its variant \emph{*ḍābba}, reconstructed forms. CDIAL calls the group non-Aryan and cites A. Master, who traced it to Dravidian (CDIAL 5539); DEDR sets it beside Dravidian words for left, such as Telugu \emph{ḍā} (DEDR 449).
\end{cousinweb}

\begin{grammarlens}[title={What you've built}]
One more word for everyday life at A2.
\end{grammarlens}

\subsection*{Your turn}
\begin{itemize}
\raggedright
  \item \emph{Say it:} \emph{ḍābũ}
  \item \emph{Say it:} \emph{ḍābũ}, once more
  \item \emph{Say it:} \emph{mandir ḍābe che} — The temple is on the left.
  \item \emph{Recall:} say the Gujarati for to waste, then the Gujarati for direction, then say \emph{ḍābũ} again
\end{itemize}

\begin{tcolorbox}[breakable,colback=teal!4,colframe=teal!35!black,title={Before you move on}]
What does \gu{ડાબું} mean? (\textquotedblleft{}Left, left-hand\textquotedblright{}.) Say it once more.
\end{tcolorbox}
\section[\texorpdfstring{\gu{જમણું} (jamṇũ)}{જમણું (jamṇũ)}]{\gu{જમણું} (jamṇũ) — right-hand, right}

\label{lesson:GU-C193-jamnu}



\begin{quote}
Before the new one: say the Gujarati for direction, then the Gujarati for left.
\end{quote}

\subsection*{You'll want to know: \gu{જમણું}}
\textbf{\gu{જમણું}} — \emph{jamṇũ} — \textquotedblleft{}right-hand, right\textquotedblright{}.

\begin{quote}
\textbf{\gu{બેંક} \gu{જમણે} \gu{છે}.} — \emph{bẽk jamṇe che} — \textquotedblleft{}The bank is on the right.\textquotedblright{}
\end{quote}

\begin{cousinweb}
\textbf{\gu{જમણું}} (\emph{jamṇũ}) is a tadbhava. CDIAL lists it, glossed \textquotedblleft{}right (not left)\textquotedblright{}, under Sanskrit \emph{jamana}, \textquotedblleft{}eating\textquotedblright{}: the right hand is the hand one eats with (CDIAL 5127).
\end{cousinweb}

\begin{grammarlens}[title={What you've built}]
One more word for everyday life at A2.
\end{grammarlens}

\subsection*{Your turn}
\begin{itemize}
\raggedright
  \item \emph{Say it:} \emph{jamṇũ}
  \item \emph{Say it:} \emph{jamṇũ}, once more
  \item \emph{Say it:} \emph{bẽk jamṇe che} — The bank is on the right.
  \item \emph{Recall:} say the Gujarati for direction, then the Gujarati for left, then say \emph{jamṇũ} again
\end{itemize}

\begin{tcolorbox}[breakable,colback=teal!4,colframe=teal!35!black,title={Before you move on}]
What does \gu{જમણું} mean? (\textquotedblleft{}Right-hand, right\textquotedblright{}.) Say it once more.
\end{tcolorbox}
\section[\texorpdfstring{\gu{સીધું} (sīdhũ)}{સીધું (sīdhũ)}]{\gu{સીધું} (sīdhũ) — straight, straight on}

\label{lesson:GU-C193-sidhu}



\begin{quote}
Before the new one: say the Gujarati for left, then the Gujarati for right-hand.
\end{quote}

\subsection*{You'll want to know: \gu{સીધું}}
\textbf{\gu{સીધું}} — \emph{sīdhũ} — \textquotedblleft{}straight, straight on\textquotedblright{}.

\begin{quote}
\textbf{\gu{સ્ટેશન} \gu{સીધું} \gu{આગળ} \gu{છે}.} — \emph{sṭeśan sīdhũ āgaḷ che} — \textquotedblleft{}The station is straight ahead.\textquotedblright{}
\end{quote}

\begin{cousinweb}
\textbf{\gu{સીધું}} (\emph{sīdhũ}) is a tadbhava. CDIAL lists it, glossed \textquotedblleft{}straightforward\textquotedblright{}, under Sanskrit \emph{siddha}, \textquotedblleft{}accomplished; prepared\textquotedblright{} (CDIAL 13401).
\end{cousinweb}

\begin{grammarlens}[title={What you've built}]
One more word for everyday life at A2.
\end{grammarlens}

\subsection*{Your turn}
\begin{itemize}
\raggedright
  \item \emph{Say it:} \emph{sīdhũ}
  \item \emph{Say it:} \emph{sīdhũ}, once more
  \item \emph{Say it:} \emph{sṭeśan sīdhũ āgaḷ che} — The station is straight ahead.
  \item \emph{Recall:} say the Gujarati for left, then the Gujarati for right-hand, then say \emph{sīdhũ} again
\end{itemize}

\begin{tcolorbox}[breakable,colback=teal!4,colframe=teal!35!black,title={Before you move on}]
What does \gu{સીધું} mean? (\textquotedblleft{}Straight, straight on\textquotedblright{}.) Say it once more.
\end{tcolorbox}
\section[\texorpdfstring{\gu{વળવું} (vaḷvũ)}{વળવું (vaḷvũ)}]{\gu{વળવું} (vaḷvũ) — to turn}

\label{lesson:GU-C193-valvu-turn}



\begin{quote}
Before the new one: say the Gujarati for right-hand, then the Gujarati for straight.
\end{quote}

\subsection*{You'll want to know: \gu{વળવું}}
\textbf{\gu{વળવું}} — \emph{vaḷvũ} — \textquotedblleft{}to turn\textquotedblright{}.

\begin{quote}
\textbf{\gu{બસ} \gu{મંદિર} \gu{પાસે} \gu{ડાબે} \gu{વળે} \gu{છે}.} — \emph{bas mandir pāse ḍābe vaḷe che} — \textquotedblleft{}The bus turns left by the temple.\textquotedblright{}
\end{quote}

\begin{cousinweb}
\textbf{\gu{વળવું}} (\emph{vaḷvũ}) is a tadbhava. It comes from Sanskrit \emph{válati}, \textquotedblleft{}turns, turns to\textquotedblright{}, and CDIAL glosses the Gujarati verb \textquotedblleft{}to turn, turn round, bend\textquotedblright{} (CDIAL 11405).
\end{cousinweb}

\begin{grammarlens}[title={What you've built}]
One more word for everyday life at A2.
\end{grammarlens}

\subsection*{Your turn}
\begin{itemize}
\raggedright
  \item \emph{Say it:} \emph{vaḷvũ}
  \item \emph{Say it:} \emph{vaḷvũ}, once more
  \item \emph{Say it:} \emph{bas mandir pāse ḍābe vaḷe che} — The bus turns left by the temple.
  \item \emph{Recall:} say the Gujarati for right-hand, then the Gujarati for straight, then say \emph{vaḷvũ} again
\end{itemize}

\begin{tcolorbox}[breakable,colback=teal!4,colframe=teal!35!black,title={Before you move on}]
What does \gu{વળવું} mean? (\textquotedblleft{}To turn\textquotedblright{}.) Say it once more.
\end{tcolorbox}
